% =====================================================================
%  THÈME 2 — Conversions en quantité de matière
%  Niveau MOYEN — Corrigé
% =====================================================================
\documentclass[11pt,a4paper]{article}
\usepackage[utf8]{inputenc}
\usepackage[T1]{fontenc}
\usepackage{lmodern}
\usepackage[french]{babel}
\usepackage{amsmath,amssymb,mathtools}
\usepackage{icomma}
\usepackage[version=4]{mhchem}
\usepackage{siunitx}
\sisetup{output-decimal-marker={,},per-mode=symbol}
\DeclareSIUnit\Molar{\mole\per\litre}
\usepackage{xcolor}
\usepackage{tikz}
\usepackage[most]{tcolorbox}
\usepackage{enumitem}
\usepackage{array}
\usepackage{fancyhdr}
\usepackage[a4paper,margin=2.1cm,headheight=26pt,headsep=8pt]{geometry}

\definecolor{primary}{RGB}{142,93,168}
\definecolor{primarydark}{RGB}{82,45,108}
\definecolor{corrcol}{RGB}{176,110,20}
\newcommand{\gmol}{\si{\gram\per\mole}}
\newcommand{\Vm}{V_{\mathrm m}}

\pagestyle{fancy}\fancyhf{}
\fancyhead[L]{\small\textcolor{primarydark}{\textbf{Fiche 5} — Thème 2 : Conversions}}
\fancyhead[R]{\small\textcolor{corrcol}{\textsc{Corrigé — Moyen}}}
\fancyfoot[C]{\small\thepage}
\renewcommand{\headrulewidth}{0.8pt}
\renewcommand{\headrule}{\hbox to\headwidth{\color{primary}\leaders\hrule height \headrulewidth\hfill}}

\newtcolorbox[auto counter]{cor}[1][]{enhanced,breakable,boxrule=1pt,arc=3pt,
  left=3mm,right=3mm,top=2.4mm,bottom=2.4mm,colback=corrcol!6,colframe=corrcol,
  fonttitle=\bfseries,coltitle=white,
  title={Corrigé~\thetcbcounter\ \ifx\relax#1\relax\relax\else\textmd{--- \itshape #1}\fi}}

\newcommand{\rep}[1]{\par\smallskip\noindent\textbf{\color{corrcol}$\Rightarrow$ #1}}

\setlength{\parindent}{0pt}
\setlength{\parskip}{3pt}

\begin{document}

\begin{center}
{\LARGE\bfseries\color{primarydark}Thème 2 — Conversions}\\[1pt]
{\large\color{corrcol}\textbf{Corrigé — Niveau Moyen}}\\
{\color{primary}\rule{0.6\linewidth}{1pt}}
\end{center}

\begin{cor}[du gramme au volume gazeux]
Mole : $n=\dfrac{m}{M}=\dfrac{16}{32}=\SI{0,50}{\mole}$.\quad
Volume : $V=n\,\Vm = 0{,}50\times22{,}4 = \boxed{\SI{11,2}{\litre}}$.
\rep{Chaîne masse $\to$ mole $\to$ volume respectée.}
\end{cor}

\begin{cor}[la densité d'un liquide]
\textbf{a)} $m=\rho\times V = 1{,}70\times100 = \SI{170}{\gram}$.\\
\textbf{b)} $n=\dfrac{m}{M}=\dfrac{170}{85}=\boxed{\SI{2,0}{\mole}}$.
\rep{La densité sert à passer du \emph{volume} à la \emph{masse} avant $n=m/M$.}
\end{cor}

\begin{cor}[peser un nombre voulu de molécules]
\textbf{a)} $n=\dfrac{N}{N_A}=\dfrac{\SI{1,20e21}{}}{\SI{6,02e23}{}}=\SI{1,99e-3}{\mole}$
($\approx\SI{2,0e-3}{\mole}$).\\
\textbf{b)} $m=n\times M = \SI{1,99e-3}{}\times182 = \SI{0,363}{\gram} = \boxed{\SI{363}{\milli\gram}}$
(soit $\approx\SI{364}{\milli\gram}$).
\rep{Chaîne entités $\to$ mole $\to$ masse. Ordre : $\div N_A$ puis $\times M$.}
\end{cor}

\begin{cor}[du comprimé à la concentration]
\textbf{a)} $\SI{755}{\milli\gram}=\SI{0,755}{\gram}$ ; $n=\dfrac{0{,}755}{151}=\SI{5,0e-3}{\mole}$.\\
\textbf{b)} $\SI{200}{\milli\litre}=\SI{0,200}{\litre}$ ;
$C=\dfrac{n}{V}=\dfrac{\SI{5,0e-3}{}}{0{,}200}=\boxed{\SI{2,5e-2}{\Molar}}$.
\rep{Deux conversions d'unités à ne pas oublier : mg $\to$ g et mL $\to$ L.}
\end{cor}

\begin{cor}[remonter à la masse molaire d'un gaz]
\textbf{a)} $n=\dfrac{V}{\Vm}=\dfrac{5{,}6}{22{,}4}=\SI{0,25}{\mole}$.\\
\textbf{b)} $M=\dfrac{m}{n}=\dfrac{5{,}0}{0{,}25}=\boxed{\SI{20}{\gram\per\mole}}$.
\rep{Le volume gazeux donne $n$ ; la masse pesée donne ensuite $M=m/n$.}
\end{cor}

\end{document}
